\section{Measure updates with the same inputs and boundaries}
The following counts run from a function's entry through its return to the supplied
continuation. They include its setup and restoring instructions, if any, and exclude
assembly, session creation, caller stores and console output. Pseudoinstructions count
as the instructions actually executed after assembly. The table update receives an
independently prepared table, so its column does not include table construction.

\begin{table}[htbp]
\centering
\caption{Retired simulated instructions for the illustrated branch-feedback implementation.}
\tagpdfsetup{table/header-rows={1}}
\begin{tabular}{lrrr}
\toprule
Dataset & Bytes & Bit update & Prepared-table update\\
\midrule
Empty & 0 & 12 & 5\\
One byte \texttt{80} & 1 & 84 & 19\\
ASCII check message & 9 & 631 & 131\\
Zero bytes & 64 & 4415 & 901\\
Octets zero through 255 & 256 & 17709 & 3589\\
Strided octets & 4096 & 282498 & 57349\\
\bottomrule
\end{tabular}
\end{table}

The last dataset has byte \(b_i=(37i+11)\bmod256\) at zero-based index \(i\).
These definitions specify the data, not just their sizes. The bit loop's conditional
feedback means equal-length messages can retire different instruction counts. The
prepared-table loop executes the same instruction sequence for each byte, but accesses
content-dependent table positions.

The separately measured table initializer retires 17420 instructions and writes
1024 table bytes. For one 256-byte message with a newly built table, construction plus
update costs \(17420+3589=21009\) instructions, exceeding the bit update's 17709.
For the 4096-byte dataset, the sum is \(17420+57349=74769\), below 282498.
The smaller update column alone does not decide whether building a table pays for a
particular amount of work.

A table can serve many updates of one file: build it once before the read loop,
then reuse it for every chunk. Rebuilding it after each read would make the buffer
size change setup cost even though the checksum remains the same. A stored constant
table is another design, with image-size and initialization tradeoffs to document.
